\section{FlashAttention：把 attention 写成 IO-aware 算法}

本节把前面所有硬件技巧汇总到 FlashAttention：KQV tiling、online softmax、forward pass 分块。

进入具体算法前，Slide 49 先总结 workload optimization 的三条主线：coalescing/fusion 减少 HBM access，tiling 把高复用数据搬到 shared memory，quantization/recomputation 则用精度或额外 compute 换 memory。FlashAttention 不是列表之外的新魔法，而是同时采用 tiling、recomputation 与 fusion，并针对 softmax 的全局归一化约束重新组织计算顺序。

\begin{figure}[H]
\centering
\includegraphics[width=0.88\textwidth]{slide-049.jpg}
\caption{Slide 49：workload performance recap；减少访问、搬到片上、以 compute/accuracy 换 memory。}
\figsource{CS336 Spring 2026 Lecture 5 官方幻灯片。}
\end{figure}

\begin{knowledgebox}{从优化清单到算法设计}
普通 kernel 优化通常保持公式和中间张量不变，只改善执行；IO-aware algorithm 会进一步问中间张量是否必须存在。Attention 的 $QK^\top$ 分数矩阵规模为 $n^2$，若完整写入 HBM 再读回做 softmax，即使 matmul 很快也会被 IO 限制。FlashAttention 的突破是证明 softmax 可以分块在线维护统计量，因此无需物化完整矩阵。
\end{knowledgebox}


\begin{importantbox}{FlashAttention 的本质}
FlashAttention 不是改变 attention 数学，而是改变计算顺序：把 Q/K/V 分块载入 SRAM，分块计算 softmax 统计量并在线更新输出，避免把完整 \(n\times n\) attention matrix 写回 HBM。
\end{importantbox}


Slide 50 展示论文中的核心技术图，并提出“dramatically accelerates attention, but how?”。正确读法是先沿数据流数 HBM transactions：Q block 留在片上，K/V blocks 分批载入，每个 tile 产生局部 logits 与 partial output；在线 softmax 用 running maximum 和 normalization sum 合并不同 tiles。算法会重算部分中间值，却省下对完整 attention matrix 的写入与读取。

\begin{figure}[H]
\centering
\includegraphics[width=0.88\textwidth]{slide-050.jpg}
\caption{Slide 50：FlashAttention 的问题设定与论文技术图。}
\figsource{CS336 Spring 2026 Lecture 5 官方幻灯片；Dao et al.}
\end{figure}

\begin{importantbox}{FlashAttention 的反直觉 tradeoff}
它可能执行比 naive attention 更多的 FLOPs，因为 backward 或 tile 边界会重算 logits；但现代 GPU 的矩阵计算远快于 HBM 搬运，所以“多算、少存”反而更快。这个结论依赖硬件比例：若序列短、kernel launch 主导、片上容量不足或实现无法使用高效 tensor-core tiles，收益会缩小，必须以目标 shape 的 benchmark 为准。
\end{importantbox}


\begin{figure}[H]
\centering
\includegraphics[width=0.88\textwidth]{slide-051.jpg}
\caption{Slide 51: Recap of attention computation}
\figsource{CS336 Spring 2026 Lecture 5 官方幻灯片。}
\end{figure}

\noindent\textbf{展开说明：}Attention computation: 3 matrix multiplies (K, Q, V) with a softmax in between

\begin{knowledgebox}{读图：Slide 51 应该怎么看}
这页是硬件机制或性能证据页。先看数据从哪里来、在哪里复用、是否写回 HBM；再判断瓶颈是 compute、memory bandwidth、thread scheduling 还是 alignment。GPU 优化不是让公式更漂亮，而是让数据在正确层级停留更久、让 tensor cores 少等数据。
\end{knowledgebox}


\begin{figure}[H]
\centering
\includegraphics[width=0.88\textwidth]{slide-052.jpg}
\caption{Slide 52: Tiling part 1: tiling for the KQV matrix multiply}
\figsource{CS336 Spring 2026 Lecture 5 官方幻灯片。}
\end{figure}

\noindent\textbf{展开说明：}This figure 1 from the paper is literally just tiling for a KQV matrix multiply..

\begin{knowledgebox}{读图：Slide 52 应该怎么看}
这页是硬件机制或性能证据页。先看数据从哪里来、在哪里复用、是否写回 HBM；再判断瓶颈是 compute、memory bandwidth、thread scheduling 还是 alignment。GPU 优化不是让公式更漂亮，而是让数据在正确层级停留更久、让 tensor cores 少等数据。
\end{knowledgebox}


\begin{figure}[H]
\centering
\includegraphics[width=0.88\textwidth]{slide-053.jpg}
\caption{Slide 53: Tiling part 2: incremental computation of the softmax}
\figsource{CS336 Spring 2026 Lecture 5 官方幻灯片。}
\end{figure}

\noindent\textbf{展开说明：}From Mikailov and Gimelshein 2018,

\begin{knowledgebox}{读图：Slide 53 应该怎么看}
这页是硬件机制或性能证据页。先看数据从哪里来、在哪里复用、是否写回 HBM；再判断瓶颈是 compute、memory bandwidth、thread scheduling 还是 alignment。GPU 优化不是让公式更漂亮，而是让数据在正确层级停留更久、让 tensor cores 少等数据。
\end{knowledgebox}


\begin{figure}[H]
\centering
\includegraphics[width=0.88\textwidth]{slide-054.jpg}
\caption{Slide 54: Putting it all together – the forward pass of flash attention}
\figsource{CS336 Spring 2026 Lecture 5 官方幻灯片。}
\end{figure}

\noindent\textbf{展开说明：}From Dao 2023, we see

\subsection{本章小结}
本节的共同问题是如何减少昂贵的数据移动并提高硬件利用率。GPU 的速度来自海量并行和专用矩阵硬件，但只有当 memory access、shape、precision 和 kernel 组织匹配时，这些速度才会真正出现。

\subsection{GPU 性能实验应该记录什么}

现在 55 张源 slide 都已纳入，最后把性能故事转成实验协议。对任意 kernel 或 IO-aware algorithm，至少需要记录输入 shape、dtype、warmup/重复次数、kernel time、端到端 time、读写 byte、实现选择和硬件型号；否则 benchmark 无法复现。若比较低精度，还要报告收敛或误差；若比较 fusion/recomputation，则要报告 peak memory；若比较 tiling/FlashAttention，则要覆盖多个 sequence length、head dimension 与 batch size，而不是只展示最好点。

\begin{longtable}{p{0.20\textwidth}p{0.32\textwidth}p{0.38\textwidth}}
\toprule
问题 & 最低证据 & 常见误判 \\
\midrule
Compute-bound 还是 memory-bound？ & Roofline 位置、achieved FLOP/s、effective bandwidth。 & 峰值 FLOP/s 很高就认为 kernel 应该快。 \\
低精度是否有效？ & 吞吐、转换开销、scale metadata、loss/误差曲线。 & 只按 bit 数推算 2 倍或 4 倍加速。 \\
Fusion 或重算是否值得？ & kernel 数、HBM byte、peak memory、额外 FLOPs。 & 局部 kernel 更快但端到端无变化。 \\
Tiling 是否匹配硬件？ & Tile shape、occupancy、shared memory、wave 数与对齐。 & 只用理论 arithmetic intensity 解释锯齿。 \\
FlashAttention 是否更快？ & 目标 batch/heads/length 下 latency、memory、numerical delta。 & 只引用论文总加速，不测自己的 shape。 \\
\bottomrule
\end{longtable}

\begin{importantbox}{最小复现实验：从 naive attention 到 IO-aware attention}
实现或调用一个会物化 $n\times n$ scores 的 attention baseline，再与 FlashAttention kernel 比较。固定 dtype 和输出容差，扫描 sequence length、head dimension 与 batch size；记录 forward/backward time、peak memory 和数值误差。随后加入一个未参与选择 tile/kernel 的 hold-out shape。若加速只出现在对齐良好的少数维度，结论应写成“特定 shape 的 kernel 优势”，而不是笼统写成 FlashAttention 永远更快。
\end{importantbox}

